%----------------------------------------------------------------------
\chapter{Software overview}
\label{ch:overview}
%----------------------------------------------------------------------

\ssat{} performs local stack consumption analysis combined with
control-flow graph traversals in order to compute worst case stack
consuption.

The compiler extracts per-subprogram stack usage and per-unit call graph
information. This information is then parsed by the call graph
generator, and the result is a complete call graph that shows the full
call tree for every entry point. This tree serves as input for the
call graph decorator that will traverse it adding the stack usage
information. Once all the information is available in the tree, the
worst case analyzer can generate an ordered list of worst case paths
regarding stack usage for every entry point in the partition. The GNAT
Programming Studio (GPS) will provide the graphical interface for
interacting with the user.

The \ssat{} software has the following components (figure
\ref{fg:overview-architecture}):

\begin{itemize}

\item Compiler

\item Call graph generator and decorator

\item Worst case analyzer

\item GNAT Programming Studio (GPS)

\end{itemize}

\begin{figure}
\centerline{\includegraphics*[scale=0.5]{overview-architecture.eps}}
\caption{Software components}
\label{fg:overview-architecture}
\end{figure}

%----------------------------------------------------------------------
\section{The compiler}
%----------------------------------------------------------------------

The compiler is in charge of generating two basic bits of information:

\begin{itemize}

\item The amount of stack used on a per-subprogram basis. The
  compiler adds the information about whether the computed amount is
  static or dynamic, and in the case of dynamic frame sizes the
  compiler indicates whether the required size can be bounded to a
  reliable value that can be used for worst case stack usage analysis.

\item The set of target subprograms that can be called from any
  source subprogram. The complete set of target subprograms cannot
  always be determined (there are indirect and dispatching calls), but
  the compiler indicates the places where these calls appear so they
  can be resolved later.

\end{itemize}

This information, which is as it is very valuable for stack usage
analysis, is part of the input for the \emph{graph generator and
  decorator} that will produce the complete stack usage tree.

%---------------------------------------------------------------------
\section{The call graph generator and decorator}
%---------------------------------------------------------------------

This component is in charge of generating the complete call graph for
every entry point in the partition, including their stack usage
information. Some entry points are automatically identified (main
entry point, task bodies, interrupt handlers). The user can also
specify any additional entry point that needs to be considered in the
analysis.

The compiler cannot always determine the required information for a
complete analysis. The complete set of target subprograms that can be
called from every point in the call tree may not be known (for
example, when there are indirect calls). Additionally, the exact
amount of stack required by every node in the call tree cannot always
be statically determined (for example, when there is dynamic stack
allocation). The \emph{graph generator and decorator} will then help
the user (or an external tool) to fill the required information; it
will generate a file containing the pattern to fill the missing pieces
of information. The user will then be able to provide more precise
information about potential target subprograms and stack usage in
\emph{user-defined call graph and stack usage files}.


%---------------------------------------------------------------------
\section{The worst case analyzer}
%---------------------------------------------------------------------

This module will traverse the decorated call graph, starting at every
entry point, and will generate a file containing an ordered list of
worst case paths regarding stack usage.

A file describing the annotated call graph (using the VCG
\refcite{VCG} graph description language) will also be generated by
this module.

%---------------------------------------------------------------------
\section{GNAT Programming Studio (GPS)}
%---------------------------------------------------------------------

GPS will be the graphical interface for interacting with the user. It
means that users will be able to provide information to the tool using
this graphical interface, and to get a graphical representation of the
call graph, stack usage, and worst case paths regarding stack usage.
